% Minimal working example: Lean code with GitHub hyperlinks
% Compile with: xelatex -shell-escape minimal-example.tex

\documentclass{article}
\usepackage{fontspec}
\setmonofont{FreeMono}  % Or any monospace font
\usepackage{hyperref}
\usepackage{minted}

% Configure minted for Lean 4
\newmintinline[lean]{lean4}{bgcolor=white}
\newminted[leancode]{lean4}{
  fontsize=\footnotesize,
  breaklines,
  breakanywhere,
  tabsize=2,
  showspaces=false
}
\usemintedstyle{tango}

% Setup for GitHub URL hyperlinks
\newcommand{\leancodeurl}{}
\newcommand{\setleancodeurl}[1]{\renewcommand{\leancodeurl}{#1}}

% Make each line number a hyperlink to GitHub
\renewcommand{\theFancyVerbLine}{\arabic{FancyVerbLine}}
\renewcommand{\FancyVerbFormatLine}[1]{%
  \href{\leancodeurl\#L\arabic{FancyVerbLine}}{#1}%
}

% Command to import code from .lean files with GitHub links
% Usage: \leancodefile[options]{local_path}{github_url}
\newcommand{\leancodefile}[3][]{%
  \setleancodeurl{#3}%
  \inputminted[fontsize=\footnotesize,breaklines,breakanywhere,tabsize=2,showspaces=false,linenos=true,#1]{lean4}{#2}%
}

\begin{document}

\title{Minimal Example: Code Import with GitHub Hyperlinks}
\author{Template Example}
\date{\today}
\maketitle

\section{Example: Importing Lean Code with GitHub Links}

Here's the \texttt{GroupAction} class definition imported from the source file.
Each line number is clickable and links to GitHub:

\leancodefile[
	firstline=20,
	lastline=23,
	firstnumber=20
]{../lean/GroupAction/Defs.lean}{https://github.com/FrankieeW/GroupAction/lean/GroupAction/Defs.lean}

\textbf{Try it:} Click any line number in the code block above to jump to that line on GitHub!

\section{How It Works}

\begin{enumerate}
	\item \texttt{\textbackslash leancodefile} imports code from a local \texttt{.lean} file
	\item \texttt{firstline/lastline} specify which lines to extract
	\item \texttt{firstnumber} sets the displayed line numbering
	\item Each line is wrapped in \texttt{\textbackslash href} linking to GitHub
	\item The GitHub URL is constructed as: \texttt{base\_url\#LN} where N is the line number
\end{enumerate}

\section{Another Example: Permutation Representation}

\leancodefile[
	firstline=66,
	lastline=69,
	firstnumber=66
]{../lean/GroupAction/Permutation.lean}{https://github.com/FrankieeW/GroupAction/lean/GroupAction/Permutation.lean}

\end{document}
